% Figure description (placeholder — figure to be generated)
% See Fig. \ref{fig:denoising-examples}
\subsection{Denoising Examples}
\label{sec:experiment-fig4}

Fig. \ref{fig:denoising-examples} presents a qualitative comparison of denoising performance on representative test signals acquired from the aluminum plate specimens described in Section~\ref{sec:experiment-setup}. Each row corresponds to a distinct method applied to the same input waveform segment, enabling direct visual comparison across processing pipelines.

The top row (left to right) shows: (a) the reference high-averaged signal obtained from 256 averages, serving as the ground-truth target; (b) the low-averaged input signal with 16 averages, exhibiting an estimated SNR of approximately 6\,dB; (c) the result of Wiener filtering; (d) the result of BM3D; (e) the result of DWT denoising; and (f) the output of the proposed PMDF-Net. Each subplot displays the time-domain waveform spanning the S0 Lamb wave arrival and the subsequent higher-order mode reflections. The bottom row presents a zoomed view of the S0 arrival region, corresponding to the dashed-boxed region indicated in the top row, to facilitate detailed assessment of waveform fidelity near the onset.

A key observation is that PMDF-Net recovers the waveform morphology of the S0 arrival with the closest agreement to the reference among all compared methods, preserving both the rising edge profile and the peak amplitude. In contrast, Wiener filtering exhibits residual noise that obscures the mode onset, BM3D shows slight amplitude attenuation, and DWT introduces visible oscillations that distort the modal structure. The zoomed S0 region (bottom row) makes these differences especially apparent. This visual comparison corroborates the quantitative SNR and NMSE results reported in Table~\ref{tab:denoising-metrics}.
